% Section X: Outlook - why TQFT is useful across QFT and condensed matter
% Purpose: Not just summary - answer "why should a QFT student care?"
%
% Answer: TQFT reorganizes gauge theory around global data, makes quantization
%         and response coefficients structurally transparent, exposes anomaly/boundary
%         relations cleanly, provides a language in which low-energy universality
%         classes and protected observables become computable.
%
% Status: NEEDS TO BE WRITTEN (write after Sections II-IX complete)
% Sources: Themes scattered throughout papers, synthesize here
%
% Writing notes:
% - This is NOT just a summary of what was covered
% - Answer: why should reader care about this reorganization?
% - What does TQFT perspective enable that ordinary QFT perspective obscures?
% - Forward-looking: where else does this framework appear?
% - ~3-5 pages

\section{Outlook: The Organizing Power of Topological Field Theory}

% TODO: Write after main content complete
% Synthesize "organizing principle" perspective from throughout paper

\subsection{What the TQFT Perspective Provides}

% TODO: Explicit enumeration of benefits
% - Quantization becomes geometric (moduli spaces, symplectic reduction)
% - Response coefficients (Hall conductivity) are topological → protected
% - Anomaly/boundary relations become structural (inflow)
% - Universality classes classified by topological data
% - Protected observables (anyonic braiding, ground state degeneracy)
% - Computable: finite-dimensional Hilbert spaces, exactly solvable

\subsection{Beyond Chern--Simons Theory}

% TODO: Brief forward look at broader landscape
% - Other TQFTs: Dijkgraaf-Witten, BF theory, state-sum models
% - Higher dimensions: 4d TQFTs, topological insulators, SPT phases
% - Topological order in condensed matter
% - String theory applications (topological string, M-theory)
% - Quantum computing with anyons
%
% Keep BRIEF - roadmap, not full development

\subsection{The Topological Bootstrap}

% TODO: Conceptual point about how topology constrains dynamics
% In TQFT, topology alone determines much of the physics
% Examples from paper:
% - Level quantization from π₃(G)
% - Hilbert space dimension from topology of Σ
% - Hall conductivity from CS level
% - Edge modes from anomaly inflow
%
% This is "topological bootstrap" - topology determines more than you might expect

\subsection{Open Questions and Future Directions}

% TODO: Honest about what's NOT covered, where to go next
% - Non-abelian quantum Hall states (ν = 5/2, Moore-Read)
% - Modular tensor categories and quantum groups (mathematical structure)
% - Relation to CFT (WZW models, Verlinde formula)
% - Higher gauge theory and higher categories
% - Topological phases of matter more broadly
%
% Point reader toward natural next readings

\subsection{Concluding Remarks}

% TODO: Bring it home
% Started: topological structure inside ordinary QFT
% Found: Chern-Simons theory where topology becomes central
% Formalized: functorial TQFT axioms
% Applied: FQHE and anyonic physics
% Clarified: what is and isn't topological
%
% Final message: TQFT is not a replacement for ordinary QFT
% It's a reorganization that makes certain phenomena transparent
% When topology controls the physics, TQFT is the right language
%
% ~1 page, conclusive but not flowery
